\chapter{Discusión}
\label{ch:discusion}

Este capítulo interpreta los resultados del Capítulo~\ref{ch:resultados} desde una perspectiva práctica. La idea central es que TRSS no reemplaza universalmente a MNE-Python; define una alternativa útil cuando la pérdida de canales vuelve insuficiente la reconstrucción puramente espacial y cuando el análisis puede asumir mayor costo computacional. La discusión se organiza alrededor de la frontera de uso entre ambos métodos, la interpretación temporal y espectral de las reconstrucciones, las implicancias de costo y las limitaciones del protocolo experimental.

\section{Lectura condicional de los resultados}

La evidencia del Capítulo~\ref{ch:resultados} muestra una ventaja descriptiva de TRSS en varias métricas de reconstrucción, especialmente MAE, RMSE, NRMSE, DTW, correlación y $R^2$. Sin embargo, esa ventaja no debe interpretarse como dominancia universal. La métrica LSD favorece a MNE en el agregado y el tiempo de cómputo también favorece a MNE de manera clara. La conclusión central, por tanto, no es que TRSS sea siempre mejor, sino que su utilidad depende del régimen de pérdida y de la prioridad del análisis.

Esta lectura condicional es coherente con la formulación metodológica del Capítulo~\ref{ch:metodologia}. El protocolo ocultó artificialmente canales presentes para poder comparar contra una verdad de terreno conocida. Cada diferencia entre métodos se calculó sobre la misma señal, la misma máscara y la misma semilla. En ese marco, la tasa de victoria y la mejora mediana son descriptores más útiles que una prueba de significancia como eje de la tesis, porque indican frecuencia y magnitud práctica de mejora bajo escenarios controlados.

\section{Por qué MNE-Python sigue siendo una referencia fuerte}

MNE-Python sigue siendo una referencia difícil porque combina un supuesto fisiológicamente razonable con una implementación madura. La interpolación por splines esféricas aprovecha la suavidad espacial inducida por la conducción volumétrica y usa de forma global los canales observados. Cuando la pérdida es baja o la señal es espacialmente suave, esta información es suficiente para reconstruir canales faltantes sin introducir una restricción temporal adicional.

Esta fortaleza explica los casos donde MNE empata o supera a TRSS. En esos escenarios, la temporalidad puede aportar poco o incluso introducir sesgo si fuerza una trayectoria demasiado regular. Por eso, la tesis no debe presentar a MNE como una línea base débil. Al contrario, su presencia fortalece la evaluación: cualquier ventaja de TRSS frente a MNE ocurre contra una referencia comunitaria realista y no contra una implementación simplificada.

\section{Aporte y límites del término temporal}

El término temporal de TRSS actúa como una restricción sobre las soluciones admisibles. Cuando la información espacial local es pobre, muchas superficies espaciales pueden ser compatibles con los canales observados. La continuidad temporal reduce esa ambigüedad al favorecer trayectorias más estables. Este mecanismo explica por qué TRSS tiende a mejorar en pérdidas cercanas o periféricas severas, donde faltan vecinos informativos alrededor del canal oculto.

El estudio de contribución temporal muestra que retirar ese término cambia el desempeño dentro de la familia TRSS. Aun así, el término temporal no debe entenderse como una fuente mágica de información ni como garantía de superioridad. Su efecto depende del patrón de pérdida, de la estructura temporal de la señal y de los hiperparámetros. Esta interpretación también explica por qué TRSS puede mejorar la forma temporal en ciertos casos y no dominar el espectro global medido por LSD.

\section{Interpretación temporal y espectral}

Las figuras de señales representativas cumplen un rol distinto al de las tablas agregadas. Las tablas muestran tendencia, mientras que las señales permiten inspeccionar plausibilidad. En los casos favorables a TRSS, la reconstrucción conserva mejor transitorios o amplitudes locales respecto de MNE. En casos favorables a MNE, la regularización espacial clásica se acerca más a la señal original. En empates prácticos, ambos métodos producen curvas similares. Esta diversidad de ejemplos evita seleccionar solo casos convenientes y muestra que la frontera de uso es real.

La comparación espectral original--MNE--TRSS aclara una tensión importante. Reducir MAE no implica preservar mejor la densidad espectral de potencia. TRSS puede acercarse más punto a punto a la señal original y aun así introducir diferencias de distribución espectral. MNE, por su prior espacial, puede conservar mejor ciertas propiedades globales del espectro. Por ello, la decisión de método depende de la aplicación posterior: análisis temporal, detección de eventos, métricas de amplitud y análisis espectral no ponderan igual los errores.

\section{Costo computacional y decisión de uso}

El costo computacional es una limitación práctica de TRSS. Aunque los tiempos por caso siguen siendo manejables en análisis fuera de línea, la diferencia frente a MNE se vuelve relevante en estudios grandes, procesamiento interactivo o búsquedas de hiperparámetros. MNE debe mantenerse como opción por defecto cuando la rapidez, la estabilidad operacional y la reproducibilidad estándar pesan más que una mejora moderada de reconstrucción.

TRSS se justifica cuando la pérdida de canales amenaza la validez de métricas posteriores y el análisis admite mayor costo. Esto ocurre especialmente en pérdidas cercanas, periféricas o severas, donde el vecindario espacial queda degradado. En esos casos, la información temporal puede estabilizar la reconstrucción y reducir errores de amplitud o forma. La frontera de decisión presentada en la Tabla~\ref{tab:ch4_decision_boundary_trss_mne} sintetiza esta recomendación: MNE para condiciones estándar; TRSS para regímenes difíciles y análisis fuera de línea.

\section{Implicancias para la práctica EEG}

Para usuarios de EEG, la recomendación práctica es mantener MNE-Python como referencia inicial. Es rápido, robusto, ampliamente usado y funciona especialmente bien cuando hay suficientes canales buenos alrededor de los faltantes. TRSS debe considerarse como herramienta complementaria en contextos donde la pérdida no es aislada o donde los segmentos temporales reconstruidos afectan directamente el análisis posterior.

Esta conclusión también orienta futuras extensiones. Si TRSS se acelera mediante matrices dispersas, factorizaciones reutilizables o criterios de parada temprana, su zona de aplicación podría ampliarse. Si además se desarrolla una regla automática de selección de hiperparámetros basada solo en información observable, se reduciría la dependencia de calibraciones costosas. Sin esas mejoras, TRSS es más adecuado para análisis planificados y trazables que para preprocesamiento rutinario de baja latencia.

\section{Limitaciones del estudio}

El protocolo usa ocultamiento artificial de canales presentes para disponer de verdad de terreno. Esta estrategia permite medir error de reconstrucción de forma directa, pero no cubre todos los problemas reales de canales EEG: saturación, ruido muscular, deriva lenta, artefactos intermitentes o contactos inestables pueden afectar a los métodos de manera distinta. Por ello, los resultados deben leerse como evidencia controlada de interpolación bajo ausencia simulada, no como validación completa frente a todo tipo de canal dañado.

La evaluación cubre varios conjuntos de datos, patrones de pérdida y métricas, pero no agota la diversidad de montajes, referencias, filtros, poblaciones o tareas. Además, las métricas reportadas evalúan calidad de señal reconstruida, no impacto directo en tareas posteriores como clasificación BCI, detección de potenciales evocados, conectividad o localización de fuentes. Un método que mejora MAE puede no mejorar necesariamente una tarea neurofisiológica posterior.

Otra limitación es la selección de hiperparámetros. El oráculo sirve como cota superior no desplegable; no representa un método realista para datos nuevos. La variante calibrada reduce parcialmente este problema, pero una aplicación práctica requiere reglas de selección basadas en propiedades observables del registro y de la máscara de canales faltantes. Finalmente, el costo de TRSS debe optimizarse antes de recomendarlo en escenarios de procesamiento masivo o tiempo real.

\section{Lecciones metodológicas}

La primera lección es que las referencias comunitarias deben evaluarse con el mismo cuidado que el método propuesto. Comparar TRSS solo contra implementaciones internas o líneas base debilitadas habría producido una conclusión artificialmente optimista. MNE-Python obliga a reconocer los escenarios donde una herramienta madura sigue siendo superior.

La segunda lección es que una evaluación multi-métrica evita conclusiones simplistas. MAE, NRMSE, DTW, LSD, coherencia y tiempo miden propiedades distintas. La tesis gana solidez precisamente porque muestra ganancias, pérdidas y empates prácticos, en vez de reducir el resultado a una única métrica o a una decisión binaria.

La tercera lección es que la trazabilidad de artefactos es parte de la contribución. Cada tabla y figura debe poder rastrearse a CSV, script y parámetros. Esto es especialmente importante en capítulos de resultados: si el lector no puede conectar una afirmación con el artefacto que la origina, la interpretación se vuelve frágil.
